\documentclass[11pt]{article}
\usepackage[a4paper,margin=1in]{geometry}
\usepackage[T1]{fontenc}
\usepackage[utf8]{inputenc}
\usepackage{lmodern}
\usepackage{amsmath,amssymb}
\usepackage{microtype}
\setlength{\parindent}{1.2em}
\setlength{\parskip}{0.25em}
\newcommand{\dd}{\,d}
\begin{document}

\begin{center}
{\Large G. Lejeune Dirichlet's Werke. II.}\par\vspace{1.2em}
{\Large XVII.}\par\vspace{0.5em}
{\LARGE \textbf{On a Theorem Concerning Series.}}\par\vspace{1em}
{\large By M. G. Lejeune Dirichlet}\par\vspace{0.5em}
{\small English translation}
\end{center}

\begin{center}
{\large \textbf{On a Theorem Concerning Series.}}
\end{center}

The theorem which I am about to establish in a new way is the same one which has served me as a lemma in several preceding memoirs. The proof which I gave of it in the first of these memoirs assumes a certain condition which, although it is satisfied in most of the applications one can make of the lemma, is not indispensable, as I have already remarked elsewhere and as will also be seen in what follows.

I begin with a very simple special case. Let \(a\) and \(b\) denote positive constants and let \(\rho\) likewise be a positive variable. Then
\[
\lim_{\rho\to0}\rho\sum_{n=0}^{n=\infty}\frac{1}{(b+na)^{1+\rho}}=\frac{1}{a},
\]
where the limit refers to the indefinite decrease of \(\rho\). To prove this special case, consider the integral
\[
\int_b^\infty\frac{dx}{x^{1+\rho}}=\frac{1}{\rho b^\rho}
\]
as the area under a curve whose abscissa is \(x\). Since this curve approaches more and more closely the \(x\)-axis, if one draws parallels to this axis through the endpoints of the ordinates corresponding to \(x=b,b+a,\) and so on, one obtains exterior and interior rectangles for the area. It follows at once that the sum whose limit is in question lies between the two quantities
\[
\frac{1}{ab^\rho}
\quad\text{and}\quad
\frac{1}{ab^\rho}+\frac{\rho}{b^{1+\rho}},
\]
whose common limit is \(1/a\).

Now to pass to the general theorem, let
\[
k_1,\quad k_2,\quad k_3,\quad \ldots,\quad k_n,\quad \ldots
\]
be positive constants arranged in increasing order of magnitude, so that \(k_n\leq k_{n+1}\). Suppose that these constants are such that, if \(t\) denotes a positive continuous variable and \(T\) the number of terms of our sequence which do not exceed \(t\), the ratio
\[
\frac{T}{t}
\]
converges to the finite limit \(\alpha\) as \(t\) increases indefinitely. Under this assumption, I say that the same quantity \(\alpha\) is also the limit of the sum
\[
\tag{1} \rho\sum_{n=1}^{n=\infty}\frac{1}{k_n^{1+\rho}}
\]
for an infinitely small value of \(\rho\).

By the assumed condition, if \(\delta\) denotes a number as small as desired, one can choose another number \(\tau\) large enough so that, whenever \(t\) exceeds \(\tau\), one always has
\[
\alpha-\delta<\frac{T}{t}<\alpha+\delta.
\]
This done, let \(N\) be the value of \(T\) corresponding to \(t=\tau\), and consider first only those terms of our series whose index \(n>N\). Such a term \(k_n\) can present two cases, according as one has
\[
k_n<k_{n+1}\quad\text{or}\quad k_n=k_{n+1}.
\]

In the first case, the value of \(T\) corresponding to \(t=k_n\) is \(n\), and one has
\[
\alpha-\delta<\frac{n}{k_n}<\alpha+\delta.
\]
In the second case, let \(k_{m'}\) be the last among the terms following \(k_n\) which is still equal to \(k_n\), and, among the preceding terms, let \(k_m\) be the first whose value is smaller than that of \(k_n\). If now we let \(t\) increase starting from \(t=k_m\), \(T\) remains constantly equal to \(m\) until \(t\) reaches the value \(k_{m'}\), and our ratio \(T/t\) thereby approaches \(m/k_{m'}\) indefinitely, while always remaining greater than \(\alpha-\delta\). Hence
\[
\frac{m}{k_{m'}}>\alpha-\delta.
\]
But by the first case we also have
\[
\frac{m'}{k_{m'}}<\alpha+\delta,
\]
and moreover
\[
m<n<m',\qquad k_n=k_{m'}.
\]
We therefore conclude, as before, that
\[
\alpha-\delta<\frac{n}{k_n}<\alpha+\delta.
\]

Putting this double inequality in the form
\[
(\alpha-\delta)^{1+\rho}\frac{\rho}{n^{1+\rho}}
<
\frac{\rho}{k_n^{1+\rho}}
<
(\alpha+\delta)^{1+\rho}\frac{\rho}{n^{1+\rho}},
\]
and summing from \(n=N+1\) to \(n=\infty\), we obtain
\[
(\alpha-\delta)^{1+\rho}\rho\sum_{n=N+1}^{n=\infty}\frac{1}{n^{1+\rho}}
<
\rho\sum_{n=N+1}^{n=\infty}\frac{1}{k_n^{1+\rho}}
<
(\alpha+\delta)^{1+\rho}\rho\sum_{n=N+1}^{n=\infty}\frac{1}{n^{1+\rho}}.
\]
But the expression \(\rho\sum 1/n^{1+\rho}\) falls under the special case examined above when one puts
\[
b=N+1,\qquad a=1,
\]
and consequently has unity as its limit. Thus the part of our sum (1) which extends from \(n=N+1\) to \(n=\infty\) will eventually remain between \(\alpha-\varepsilon\) and \(\alpha+\varepsilon\), where \(\varepsilon\) is only slightly larger than \(\delta\) and therefore, like \(\delta\), arbitrarily small. At the same time the first part, which has only the finite number \(N\) of terms, plainly converges to zero. It follows that the limit of expression (1) is indeed \(\alpha\), Q.E.D.
\end{document}
